% ============================================================
% chapter2.tex  —  Chapitre 2 : Development Methodology
% ============================================================

\chapter{Development Methodology}


\noindent
Having identified the gaps that justify this work,
the next logical question is how the proposed solution
was conceived, structured, and built.
This chapter describes the development methodology adopted
for the DecisioBI platform — from the overall architectural
choices to the modular decomposition, the technology stack,
the data processing pipeline, and the testing strategy.
The goal is not to catalogue every library or configuration file,
but to explain the \emph{reasoning} behind each major technical decision
and to show that the implementation follows a coherent,
repeatable, and verifiable process.


% ============================================================
\section{General Approach: Iterative and Modular Development}
% ============================================================

\noindent
The development followed an \textbf{iterative and incremental}
approach. Rather than attempting to design the entire system
upfront and implement it in a single pass, the project was
decomposed into numbered modules (M1 through M9),
each addressing a specific functional domain.
Each module was developed, tested, and documented independently
before being integrated into the whole.

\vspace{0.4cm}
\noindent
This strategy was chosen for three practical reasons.
First, it allowed early validation of core assumptions —
particularly around data ingestion and cleaning —
before investing effort in downstream features.
Second, it reduced integration risk: each module
exposes a well-defined REST API, and modules communicate
through direct API calls and Celery tasks rather than
tight code coupling. Third, it mirrors the reality
of a thesis project with constrained timelines:
delivering a working M1 + M2 + M3 pipeline
is more valuable than delivering a half-finished
system that attempts everything at once.

\vspace{0.4cm}
\noindent
The modules are organised in six phases:

\begin{enumerate}
  \item \textbf{Phase 1 — Foundation:} Authentication and authorisation (M9),
        providing the security layer that all subsequent modules depend on.
  \item \textbf{Phase 2 — Data Acquisition:} Data ingestion (M1),
        handling file upload, preview, ERP template mapping,
        and async processing.
  \item \textbf{Phase 3 — Data Quality:} Data cleaning (M2)
        and inter-source conflict resolution (M3),
        ensuring that data entering the analytical layer
        is reliable and coherent.
  \item \textbf{Phase 4 — Analytics:} KPI calculation (M4)
        and dashboard construction (M5),
        transforming clean data into decision-relevant indicators.
  \item \textbf{Phase 5 — Intelligence:} AI-powered interpretation,
        adding a layer of automated insight.
  \item \textbf{Phase 6 — Interaction:} Natural-language chatbot
        and notification system, making the platform
        accessible to non-technical users.
\end{enumerate}


% ============================================================
\section{Architectural Design}
% ============================================================

\noindent
The platform adopts a \textbf{client–server architecture}
with a clear separation between the presentation layer
(frontend) and the business logic layer (backend).
This separation is not merely aesthetic — it reflects
a deliberate choice to keep each layer independently
testable, deployable, and replaceable.

% --- 2.2.1 Backend ---
\subsection{Backend: Django REST Framework}

\noindent
The backend is built on \textbf{Django 6.0.3}
with the \textbf{Django REST Framework (DRF) 3.17.0}.
Django was chosen for its mature ecosystem,
its built-in ORM that abstracts PostgreSQL interactions,
its signal mechanism that enables event-driven workflows,
and its extensive documentation — a non-trivial advantage
when development is carried out by a small team.

\vspace{0.4cm}
\noindent
Each functional domain is implemented as a separate
Django application within a single project:

\begin{center}
\begin{tabular}{|l|l|l|}
\hline
\textbf{App} & \textbf{Module} & \textbf{Responsibility} \\
\hline
\texttt{authentication} & M9 & JWT auth, RBAC, user management \\
\texttt{ingestion} & M1 & File upload, preview, ERP templates \\
\texttt{nettoyage} & M2 & Cleaning rules, pipelines, quality scoring \\
\texttt{conflits} & M3 & Conflict detection, resolution strategies \\
\texttt{kpi} & M4/M5 & KPI calculation, forecasting, pivots \\
\texttt{dashboard} & M5 & Preferences, saved views, widgets \\
\texttt{ia\_interpretation} & M6 & AI analysis with Groq API \\
\texttt{chatbot} & M8 & Multi-turn conversations, NL2SQL \\
\hline
\end{tabular}
\end{center}

\vspace{0.4cm}
\noindent
Each app follows a consistent internal structure:
\texttt{models.py} defines the data schema,
\texttt{serializers.py} handles validation and serialisation,
\texttt{views.py} exposes the API logic,
\texttt{services.py} encapsulates business rules
(separated from the HTTP layer for testability),
and \texttt{urls.py} maps endpoints.
This consistency is not accidental — it is a deliberate
convention that reduces cognitive overhead
when navigating between modules.

% --- 2.2.2 Frontend ---
\subsection{Frontend: React with TanStack Router}

\noindent
The frontend is a single-page application built with
\textbf{React 19}, \textbf{TypeScript 5.8},
and \textbf{TanStack Router} for file-based routing.
The choice of React reflects its dominance in the
ecosystem and the availability of mature component libraries.
TypeScript was adopted to catch type-related errors
at compile time rather than at runtime —
a significant quality gain for a project
where the frontend interacts with a complex REST API.

\vspace{0.4cm}
\noindent
The UI layer uses \textbf{shadcn/ui} (New York style)
on top of \textbf{Tailwind CSS 4.2},
providing accessible, composable components
without imposing a heavy runtime framework.
Charts are rendered with \textbf{Recharts 2.15},
chosen for its declarative API and React-native integration.

\vspace{0.4cm}
\noindent
The frontend communicates with the backend
through a singleton \texttt{ApiClient} class
that centralises authentication headers,
error handling, and token refresh logic.
All API calls are typed — the frontend defines
TypeScript interfaces that mirror the backend serializers,
ensuring that type mismatches between layers
are caught during development rather than in production.

% --- 2.2.3 Database ---
\subsection{Database: PostgreSQL}

\noindent
\textbf{PostgreSQL 18.1} was selected as the primary data store.
Its support for JSON fields, partial indexes, and
complex queries makes it well-suited for a BI platform
where data shapes vary across sources.
The schema comprises 33 tables across 9 modules,
with approximately 400 columns, 50 foreign key relationships,
and 30 indexes.

\vspace{0.4cm}
\noindent
Redis serves as the message broker for Celery
and as a caching layer for frequently accessed data.
This separation ensures that long-running tasks
(ingestion, cleaning)
do not block the HTTP request–response cycle.

% --- 2.2.4 Asynchronous Processing ---
\subsection{Asynchronous Processing: Celery}

\noindent
Several operations — file ingestion, data cleaning,
conflict detection, KPI calculation —
are inherently time-consuming.
They are delegated to \textbf{Celery 5.6.3} workers
using Redis as the broker.
This ensures that the API remains responsive
while background tasks execute,
and that task progress can be tracked
through status updates and notifications.

\vspace{0.4cm}
\noindent
The workflow follows a consistent pattern:
the API enqueues a Celery task and returns a job ID
immediately; the frontend polls the job endpoint
for status updates; when the task completes,
the user is notified and can proceed to the next step.
Each stage requires explicit user confirmation
before advancing — for example, the user must
validate cleaning results before they enter
the conflict detection and KPI calculation stages.


% ============================================================
\section{Technology Stack Summary}
% ============================================================

\noindent
Table~\ref{tab:stack} summarises the main technologies
and their roles in the platform.

\begin{table}[h!]
\centering
\caption{Technology stack of the DecisioBI platform.}
\label{tab:stack}
\begin{tabular}{|l|l|l|}
\hline
\textbf{Layer} & \textbf{Technology} & \textbf{Version} \\
\hline
Backend framework & Django + DRF & 6.0.3 / 3.17.0 \\
Language (backend) & Python & 3.12.x \\
Frontend framework & React & 19.2.0 \\
Language (frontend) & TypeScript & 5.8.3 \\
Build tool & Vite & 7.3.1 \\
Router & TanStack Router & 1.168.25 \\
UI components & shadcn/ui + Tailwind CSS & New York / 4.2.1 \\
Charts & Recharts & 2.15.4 \\
Database & PostgreSQL & 18.1 \\
Cache \& broker & Redis & 7.4.0 \\
Task queue & Celery & 5.6.3 \\
Authentication & SimpleJWT & 5.5.1 \\
AI model & Groq API (Llama 3.3 70B) & -- \\
Deployment (frontend) & Cloudflare Workers & via wrangler \\
\hline
\end{tabular}
\end{table}


% ============================================================
\section{Modular Decomposition and Data Flow}
% ============================================================

\noindent
The platform's processing pipeline follows a linear
but event-driven sequence, illustrated
in Figure~\ref{fig:pipeline}.

\begin{figure}[h!]
\centering
\fbox{\parbox{0.9\textwidth}{\centering
\vspace{0.8cm}
\textbf{Upload (M1)} $\rightarrow$ \textbf{Cleaning (M2)} $\rightarrow$ \textbf{Conflict Detection (M3)} $\rightarrow$ \textbf{KPI Calculation (M4)} $\rightarrow$ \textbf{Dashboard (M5)} $\rightarrow$ \textbf{AI Analysis (M6)} $\rightarrow$ \textbf{Chatbot (M8)}
\vspace{0.8cm}
}}
\caption{Simplified data processing pipeline.}
\label{fig:pipeline}
\end{figure}

\noindent
The pipeline is semi-automatic.
After uploading a file, the system automatically
performs ingestion and schema profiling.
The user is then notified and can trigger
the cleaning stage, where the system analyses
the data and proposes a set of automated rules.
The user reviews, adjusts, and validates these rules
before they are applied — a deliberate design choice
that balances automation with control.
After cleaning, the user can run conflict detection
and KPI calculation.
Notifications are generated at each stage
(file loaded, cleaning started, cleaning completed,
conflicts detected, KPIs calculated),
keeping the user informed without requiring them
to monitor the system actively.

\vspace{0.4cm}
\noindent
This design has two important consequences.
First, it ensures data consistency:
KPIs are never computed on uncleaned data,
and conflicts are detected before they can
pollute downstream analytics.
Second, it makes the system auditable:
every transformation is logged, and the user
can trace the history of any data source
from upload through to the final dashboard.


% ============================================================
\section{Data Ingestion (M1)}
% ============================================================

\noindent
The ingestion module is the entry point
of the entire pipeline.
It accepts three file formats — \textbf{CSV},
\textbf{Excel} (.xlsx, .xls), and \textbf{JSON} —
and performs the following steps:

\begin{enumerate}
  \item \textbf{File validation:} format check, size limits,
        duplicate detection via MD5 checksum.
  \item \textbf{Schema profiling:} automatic column type inference,
        null rate calculation, unique value counting.
  \item \textbf{Preview:} the first rows are returned to the user
        for visual confirmation before committing.
  \item \textbf{ERP templates:} predefined column mappings
        for common export formats (sales, inventory,
        customers, suppliers, general ledger),
        reducing manual mapping effort.
  \item \textbf{Async ingestion:} the file is processed
        by a Celery task that stores raw data
        in the \texttt{RawData} table,
        linked to a \texttt{DataSource} record
        that tracks status, schema, and provenance.
\end{enumerate}

\noindent
RBAC is enforced at this stage:
only users with the \textit{admin} or \textit{analyst}
role can upload data; \textit{viewer} accounts
receive a 403 Forbidden response.


% ============================================================
\section{Data Cleaning (M2)}
% ============================================================

\noindent
The cleaning module is the most granular
component of the platform.
It provides a rule engine with 18+ cleaning rule types,
including:

\begin{itemize}
  \item Structural operations: remove empty rows,
        rename columns, split/merge columns.
  \item Value transformations: standardise text,
        fill missing values (mean, median, mode, constant),
        regex replace, value mapping.
  \item Format validation: date parsing,
        amount normalisation, email/phone format checks.
  \item Quality scoring: a composite score
        computed from completeness, consistency,
        and accuracy metrics.
\end{itemize}

\noindent
Rules are grouped into \textbf{pipelines} —
ordered sequences of rules applied in sequence.
The user can preview the effect of a pipeline
before applying it, receiving a diff sample
and a quality score comparison.
After application, a \textbf{before/after comparison}
is generated, and the user must explicitly
validate or reject the cleaning results.
This two-step process (preview then validate)
is a deliberate design choice:
it gives the user control over data transformations
while automating the tedious parts.


% ============================================================
\section{Conflict Resolution (M3)}
% ============================================================

\noindent
After cleaning, the system automatically
scans for inter-source conflicts —
inconsistencies between different data sources
that cover overlapping entities or time periods.
Four conflict types are detected:

\begin{enumerate}
  \item Duplicate records (hash-based detection).
  \item Missing values across sources.
  \item Data type inconsistencies.
  \item Format inconsistencies (regex-based).
\end{enumerate}

\noindent
The \texttt{ConflictType} model is extensible —
users can create additional custom types
via the API, but the built-in detection engine
covers these four core categories.

\noindent
For each conflict, the system recommends
a resolution strategy from seven options:
\textit{manual override}, \textit{auto-merge},
\textit{default value}, \textit{user selected},
\textit{majority vote}, \textit{latest value},
and \textit{discard}.
The user can accept the recommendation
or override it.
Every resolution is logged in an activity trail,
providing full auditability.


% ============================================================
\section{KPI Calculation and Analytics (M4/M5)}
% ============================================================

\noindent
KPIs are the analytical core of the platform.
Each KPI is defined by a formula —
which can be SQL, Python, or Excel-style —
and is automatically calculated
whenever the underlying data changes.
This \textbf{event-driven calculation}
ensures that indicators are always up to date
without requiring manual recalculation.

\vspace{0.4cm}
\noindent
The KPI module provides three analytical capabilities:

\begin{itemize}
  \item \textbf{Variance analysis:} comparison against
        targets and previous periods.
  \item \textbf{Anomaly detection:} Z-score analysis
        flagging values that deviate significantly
        from historical patterns.
  \item \textbf{Forecasting:} linear regression
        with 95\% confidence intervals,
        trend direction, and R-squared goodness of fit.
\end{itemize}

\noindent
The M5 extension adds configurable pivot tables
with eight aggregation functions,
intelligent field aliasing with bilingual support
(French/English), persistent user preferences,
and saved dashboard views.


% ============================================================
\section{Artificial Intelligence Integration}
% ============================================================

\noindent
AI is integrated at two levels
of the platform.

\vspace{0.4cm}
\noindent
\textbf{KPI interpretation} uses the Groq API
with the Llama 3.3 70B model to generate
natural-language analyses of KPI data.
The user selects a data source and specific indicators,
and the model produces a structured interpretation
including trend analysis, comparisons, and recommendations.
This feature is designed for decision-makers
who need quick insights without navigating
complex dashboards.

\vspace{0.4cm}
\noindent
\textbf{The chatbot} provides a conversational
interface to the platform.
It supports multi-turn conversations,
intent recognition, entity extraction,
and natural-language-to-SQL (NL2SQL) generation.
The user can ask questions in plain French
and receive data-driven answers,
lowering the barrier to adoption
for non-technical staff.


% ============================================================
\section{Security and Access Control}
% ============================================================

\noindent
Security is implemented at multiple layers.

\vspace{0.4cm}
\noindent
\textbf{Authentication} uses JSON Web Tokens (JWT)
with access tokens valid for 60 minutes
and refresh tokens valid for 7 days.
Token rotation and blacklisting are enabled,
ensuring that revoked sessions cannot be reused.

\vspace{0.4cm}
\noindent
\textbf{Authorisation} follows a Role-Based Access
Control (RBAC) model with three roles:
\textit{admin} (full access),
\textit{analyst} (read and write data, no user management),
and \textit{viewer} (read-only access).
Each API endpoint checks the caller's role
before executing the request.

\vspace{0.4cm}
\noindent
\textbf{Data isolation} is enforced at the database level:
each organisation's data is scoped
through foreign key relationships,
preventing cross-tenant data leakage.
SQL injection is prevented through Django's ORM
and parameterised queries.
Python formula execution is sandboxed
to prevent arbitrary code execution.
Password storage uses PBKDF2 hashing (Django's default),
and API tokens are stored as SHA-256 hashes.


% ============================================================
\section{Testing Strategy}
% ============================================================

\noindent
Testing follows the \textbf{Arrange–Act–Assert} pattern
and is organised per module in three categories:

\begin{enumerate}
  \item \textbf{Model tests:} field constraints,
        relationships, computed fields, timestamps.
  \item \textbf{API tests:} status codes,
        response structure, filtering, pagination,
        permissions.
  \item \textbf{Security tests:} authentication requirements,
        authorisation enforcement, data isolation,
        input validation (SQL injection, XSS,
        malicious inputs), rate limiting.
\end{enumerate}

\noindent
The test framework is \textbf{pytest 9.0.2}
with \textbf{pytest-django} for database access
and \textbf{pytest-cov} for coverage measurement.
Custom test fixtures are used instead of
external factory libraries,
keeping the dependency footprint minimal.

\vspace{0.4cm}
\noindent
The coverage target is \textbf{70\%} minimum,
configured as a hard failure threshold
in \texttt{pytest.ini}.
Current results show 161+ tests across all modules,
with a 94.8\% pass rate (128 out of 135 tests passing).
The M4 KPI module alone has 109 tests,
covering models, services, serializers,
integration scenarios, and security.

\vspace{0.4cm}
\noindent
This testing discipline serves two purposes
in the context of this thesis.
First, it provides empirical evidence
that the platform works as specified —
each module can be verified independently.
Second, it documents the expected behaviour
of the system: reading the test suite
is in many cases clearer than reading
the implementation code,
making the tests a form of executable specification.


% ============================================================
\section{Deployment Architecture}
% ============================================================

\noindent
The backend runs on a server with Python 3.12,
PostgreSQL 18.1, and Redis 7.4.
Celery workers run as background processes,
consuming tasks from Redis queues.
The frontend is deployed on \textbf{Cloudflare Workers},
a serverless edge platform that provides
low-latency access from any geographic location —
a relevant consideration for users
in West Africa where centralised hosting
can introduce significant latency.

\vspace{0.4cm}
\noindent
In development, the Vite dev server proxies
\texttt{/api} requests to the Django backend
running on \texttt{localhost:8000}.
In production, the frontend is built as static assets
served by Cloudflare, with API calls directed
to the backend's public URL
configured via the \texttt{VITE\_API\_BASE\_URL}
environment variable.


% ============================================================
\section{Chapter Summary}
% ============================================================

\noindent
The methodology described in this chapter
is characterised by three principles.
The first is \textbf{modularity}: each functional domain
is encapsulated in a self-contained application
with its own models, services, and API,
reducing coupling and enabling independent testing.
The second is \textbf{guided automation}:
the pipeline automates tedious tasks
while keeping the user in control at critical
decision points, ensuring data quality
through deliberate validation steps.
The third is \textbf{verifiability}:
every module is accompanied by a comprehensive test suite
that documents its expected behaviour
and provides empirical evidence of correctness.

\vspace{0.4cm}
\noindent
These principles are not abstract preferences —
they are responses to the specific constraints
identified in Chapter~1: small teams, manual processes,
limited resources, and the need for tools
that work reliably without dedicated IT support.
The next chapter presents the platform itself,
showing how these principles translate
into a functional system.

\vspace{0.7cm}
\begin{center}
  \textcolor{BITred!40}{\rule{0.35\linewidth}{0.5pt}}
  \hspace{0.5em}
  \textcolor{BITred}{$\star$}
  \hspace{0.5em}
  \textcolor{BITred!40}{\rule{0.35\linewidth}{0.5pt}}
\end{center}
